%! AP Statistics — Unit 1, Lesson 1.4 — Slides (Beamer)
\documentclass[aspectratio=169, 11pt]{beamer}
\usepackage{apstats-beamer}

%=============================================================================
\begin{document}

% ─────────────────────────────────────────────────────────────────────────────
% SLIDE 1 — LESSON TITLE
% ─────────────────────────────────────────────────────────────────────────────
{
\setbeamercolor{background canvas}{bg=navy}
\begin{frame}[plain]
  \begin{tikzpicture}[remember picture, overlay]
    \fill[navylight] (current page.north west)
      rectangle ([yshift=-1.35cm]current page.north east);
  \end{tikzpicture}
  \vspace{2.8cm}
  \hspace{0.55cm}%
  \begin{minipage}{0.9\textwidth}
    {\color{goldacc}\bfseries\small LESSON 1.4}\\[0.5cm]
    {\color{white}\bfseries\LARGE Representing a Categorical Variable\\[0.25cm]
      with Graphs}\\[1.4cm]
    {\color{skymid}\small AP Statistics~$\cdot$~Unit 1: Exploring One-Variable Data}
  \end{minipage}
\end{frame}
}

% ─────────────────────────────────────────────────────────────────────────────
% SLIDE 2 — WARM-UP
% ─────────────────────────────────────────────────────────────────────────────
{
\setbeamercolor{background canvas}{bg=sky}
\begin{frame}[t, plain]
  \navyheader{Warm-Up}
  \vspace{1em}%
  \hspace{0.55cm}%
  \begin{minipage}{0.9\textwidth}
    \small

    Recall the transportation survey from Lesson 1.3 ($n = 25$):
    Car~40\%, Bus~32\%, Walk~16\%, Bike~12\%.

    \vspace{0.5em}
    \begin{enumerate}
      \setlength{\itemsep}{0.7em}\setlength{\topsep}{0em}
      \item What is the most common mode of transportation and what percentage
            of students use it?
      \item If you wanted to show this distribution \textit{visually} --- so a
            reader could compare categories at a glance --- what type of picture
            would you draw?
      \item What information does a visual display show that the table does not
            show as easily? What does the table show that a visual might not?
    \end{enumerate}

    \vspace{0.5cm}
    {\color{navy}\itshape\small
      Today we formalize two visual displays --- bar charts and pie charts ---
      and learn how graphs can mislead even with correct data.}
  \end{minipage}
\end{frame}
}

% ─────────────────────────────────────────────────────────────────────────────
% SLIDE 3 — PRIMARY OBJECTIVE
% ─────────────────────────────────────────────────────────────────────────────
{
\setbeamercolor{background canvas}{bg=sky}
\begin{frame}[t, plain]
  \begin{tikzpicture}[remember picture, overlay]
    \fill[navy] (current page.north west)
      rectangle ([yshift=-0.25cm]current page.north east);
  \end{tikzpicture}
  \vspace{0.45cm}\par
  \hspace{0.55cm}%
  \begin{minipage}{0.9\textwidth}

    \sectionlabel{Primary Objective}

    \normalsize
    Students will construct and interpret \textbf{bar charts} and
    \textbf{pie charts} for categorical variables, and will identify
    \textbf{misleading features} in graphs (Big Idea: UNC).

    \vspace{0.45cm}

    \normalsize\textbf{\color{navy}By the end of this lesson, I will be able to:}

    \smallskip
    \begin{itemize}
      \setlength{\itemsep}{0.3em}\setlength{\topsep}{0.2em}
      \item Construct a correctly labeled bar chart (frequency or relative frequency).
      \item Compute sector angles and construct a pie chart.
      \item Interpret bar and pie charts in context using AP-quality language.
      \item Identify and explain misleading features of a graph.
    \end{itemize}

  \end{minipage}
\end{frame}
}

% ─────────────────────────────────────────────────────────────────────────────
% SLIDE 4 — HOOK: BAD GRAPHS
% ─────────────────────────────────────────────────────────────────────────────
{
\setbeamercolor{background canvas}{bg=hookbg}
\begin{frame}[t, plain]
  \begin{tikzpicture}[remember picture, overlay]
    \fill[goldacc] (current page.north west)
      rectangle ([yshift=-1.35cm]current page.north east);
    \node[anchor=west, text=white, font=\bfseries\large,
          inner xsep=0.55cm, inner ysep=0pt]
      at ([yshift=-0.68cm]current page.north west)
      {Hook --- Can a Graph Lie?};
  \end{tikzpicture}
  \vspace{1.2cm}\par
  \vspace{1em}%
  \hspace{0.55cm}%
  \begin{minipage}{0.9\textwidth}
    \small

    % ── Two bar charts drawn with TikZ ───────────────────────────────────────
    % Chart A: y-axis 950–1010  →  X bar height 0.833, Y bar height 2.083
    % Chart B: y-axis 0–1050   →  X bar height 2.310, Y bar height 2.381
    \begin{tikzpicture}[x=1.22cm, y=1.22cm]

      % ── Chart A (left) ── y starts at 950 ──────────────────────────────────
      \node[white, font=\small\bfseries, anchor=south] at (1.5, 2.5) {Graph A};
      \node[white, font=\scriptsize, anchor=north] at (1.5, 2.5)
            {(y-axis starts at 950)};
      \draw[white, thick] (0, 2.5) -- (0, 0) -- (3.0, 0);
      \node[white, left, font=\scriptsize] at (0, 0)   {950};
      \node[white, left, font=\scriptsize] at (0, 2.5) {1010};
      % bars
      \fill[goldacc!85]  (0.35, 0) rectangle (1.15, 0.833);
      \fill[skymid!85]   (1.65, 0) rectangle (2.45, 2.083);
      % x-labels
      \node[white, below, font=\small] at (0.75, 0) {X};
      \node[white, below, font=\small] at (2.05, 0) {Y};
      % value labels
      \node[white, above, font=\scriptsize] at (0.75, 0.833) {970};
      \node[white, above, font=\scriptsize] at (2.05, 2.083) {1000};

      % ── Chart B (right) ── y starts at 0 ───────────────────────────────────
      \begin{scope}[xshift=4.4cm]
        \node[white, font=\small\bfseries, anchor=south] at (1.5, 2.5) {Graph B};
        \node[white, font=\scriptsize, anchor=north] at (1.5, 2.5)
              {(y-axis starts at 0)};
        \draw[white, thick] (0, 2.5) -- (0, 0) -- (3.0, 0);
        \node[white, left, font=\scriptsize] at (0, 0)   {0};
        \node[white, left, font=\scriptsize] at (0, 2.5) {1050};
        % bars: X=970/1050*2.5=2.310, Y=1000/1050*2.5=2.381
        \fill[goldacc!85]  (0.35, 0) rectangle (1.15, 2.310);
        \fill[skymid!85]   (1.65, 0) rectangle (2.45, 2.381);
        \node[white, below, font=\small] at (0.75, 0) {X};
        \node[white, below, font=\small] at (2.05, 0) {Y};
        \node[white, above, font=\scriptsize] at (0.75, 2.310) {970};
        \node[white, above, font=\scriptsize] at (2.05, 2.381) {1000};
      \end{scope}

    \end{tikzpicture}

    \vspace{0.35cm}
    {\color{white}\small
      Graph A: Y looks \textbf{2.5$\times$ taller} than X. \quad
      Graph B: the bars are \textbf{nearly identical}.\\[5pt]
      \textit{Same data. Non-zero y-axis. Very different story.}}
  \end{minipage}
\end{frame}
}

% ─────────────────────────────────────────────────────────────────────────────
% SLIDE 5 — VOCABULARY
% ─────────────────────────────────────────────────────────────────────────────
{
\setbeamercolor{background canvas}{bg=greenbg}
\begin{frame}[t, plain]
  \navyheader{Vocabulary}
  \vspace{1em}%
  \hspace{0.55cm}%
  \begin{minipage}{0.9\textwidth}
    \small
    \begin{description}[leftmargin=0em, itemsep=0.45em, font=\bfseries\color{navy}]
      \item[Bar chart] A graph of separated bars; heights represent frequency or
            relative frequency for each category. Bars do \textbf{not} touch.
      \item[Relative frequency bar chart] Bar chart with relative frequency on the
            vertical axis; preferred for comparing groups.
      \item[Pie chart] A circle divided into sectors proportional to relative
            frequencies; best for part-to-whole relationships with few categories.
      \item[Segmented bar chart] Bars divided into segments to show composition;
            used when comparing distributions across groups.
      \item[Misleading graph] A visual that distorts differences through a non-zero
            axis, 3-D effects, or inconsistent scale.
    \end{description}
  \end{minipage}
\end{frame}
}

% ─────────────────────────────────────────────────────────────────────────────
% SLIDE 6 — BUILDING A BAR CHART
% ─────────────────────────────────────────────────────────────────────────────
{
\setbeamercolor{background canvas}{bg=white}
\begin{frame}[t, plain]
  \navyheader{Building a Bar Chart --- Four Requirements}
  \vspace{1em}%
  \hspace{0.55cm}%
  \begin{minipage}{0.9\textwidth}
    \small

    \begin{enumerate}
      \setlength{\itemsep}{0.55em}
      \item \textbf{Title} --- states what the graph shows (variable and group).
      \item \textbf{Labeled horizontal axis} --- lists each category.
      \item \textbf{Labeled vertical axis} --- frequency \textit{or} relative
            frequency; scale \textbf{starts at 0}.
      \item \textbf{Gaps between bars} --- signals that the variable is categorical
            (not continuous).
    \end{enumerate}

    \vspace{0.5cm}
    {\color{navy}\bfseries\small Frequency vs.\ Relative Frequency:}\\[6pt]
    \small
    \begin{tabular}{@{} l l @{}}
      \textbf{Frequency bar chart} & y-axis shows \textbf{counts}; use for a single group.\\[4pt]
      \textbf{Relative frequency bar chart} & y-axis shows \textbf{proportions/\%};
        use to compare groups.
    \end{tabular}

    \vspace{0.4cm}
    {\color{slate}\footnotesize
      AP exam: an unlabeled graph earns partial or no credit even if the bar heights are correct.}
  \end{minipage}
\end{frame}
}

% ─────────────────────────────────────────────────────────────────────────────
% SLIDE 7 — WORKED EXAMPLE: BAR CHART
% ─────────────────────────────────────────────────────────────────────────────
{
\setbeamercolor{background canvas}{bg=white}
\begin{frame}[t, plain]
  \navyheader{Worked Example --- Pet Ownership ($n = 40$)}
  \vspace{1em}%
  \hspace{0.55cm}%
  \begin{minipage}{0.9\textwidth}
    \small

    \renewcommand{\arraystretch}{1.5}
    \begin{tabular}{@{} >{\bfseries}p{0.16\textwidth}
                        >{\centering\arraybackslash}p{0.14\textwidth}
                        >{\centering\arraybackslash}p{0.20\textwidth} @{}}
      \rowcolor{sky}
      \textbf{Pet type} & \textbf{Freq.} & \textbf{Rel.\ Freq.} \\
      \midrule
      Dog   & 12 & 0.30 \\
      Cat   & 8  & 0.20 \\
      Fish  & 6  & 0.15 \\
      Bird  & 4  & 0.10 \\
      None  & 10 & 0.25 \\
      \midrule
      \rowcolor{sky}\textbf{Total} & 40 & 1.00 \\
    \end{tabular}

    \vspace{0.4cm}
    {\color{navy}\bfseries\small AP-style interpretation:}\\[4pt]
    \small\textit{``Of the 40 households surveyed, dog was the most common pet
    type (30\%). Bird was the least common (10\%). The distribution is fairly
    spread across categories, with no single type dominating.''}

    \vspace{0.3cm}
    {\color{slate}\footnotesize
      \textbf{Common error:} making bars touch. Remember --- categorical $\to$ gaps.}
  \end{minipage}
\end{frame}
}

% ─────────────────────────────────────────────────────────────────────────────
% SLIDE 8 — PIE CHART
% ─────────────────────────────────────────────────────────────────────────────
{
\setbeamercolor{background canvas}{bg=sky}
\begin{frame}[t, plain]
  \navyheader{Constructing a Pie Chart}
  \vspace{1em}%
  \hspace{0.55cm}%
  \begin{minipage}{0.9\textwidth}
    \small

    \textbf{\color{navy}Sector angle formula:}
    \[
      \text{Sector angle} = \text{relative frequency} \times 360^\circ
    \]

    \vspace{0.35cm}
    \textbf{\color{navy}Pet ownership --- angles:}

    \vspace{0.3cm}
    \renewcommand{\arraystretch}{1.4}
    \begin{tabular}{@{} >{\bfseries}p{0.14\textwidth}
                        >{\centering\arraybackslash}p{0.18\textwidth}
                        >{\centering\arraybackslash}p{0.20\textwidth} @{}}
      \rowcolor{navy}
      {\color{white}Category} & {\color{white}Rel.\ Freq.} & {\color{white}Angle} \\
      Dog  & 0.30 & $0.30 \times 360 = 108^\circ$ \\
      Cat  & 0.20 & $0.20 \times 360 = 72^\circ$  \\
      Fish & 0.15 & $0.15 \times 360 = 54^\circ$  \\
      Bird & 0.10 & $0.10 \times 360 = 36^\circ$  \\
      None & 0.25 & $0.25 \times 360 = 90^\circ$  \\
    \end{tabular}

    \vspace{0.4cm}
    {\color{navy}\small\bfseries Bar vs.\ Pie:}
    Bar charts are easier to compare; pie charts emphasize part-to-whole.
    Prefer bar charts when there are more than 5--6 categories.
  \end{minipage}
\end{frame}
}

% ─────────────────────────────────────────────────────────────────────────────
% SLIDE 9 — MISLEADING GRAPHS
% ─────────────────────────────────────────────────────────────────────────────
{
\setbeamercolor{background canvas}{bg=sky}
\begin{frame}[t, plain]
  \navyheader{Misleading Graphs --- Three Common Tricks}
  \vspace{1em}%
  \hspace{0.55cm}%
  \begin{minipage}{0.9\textwidth}
    \small

    \begin{description}[leftmargin=0em, itemsep=0.5em, font=\bfseries\color{navy}]
      \item[Trick 1 --- Non-zero y-axis]
            The axis starts above 0, making small differences appear large.\\
            \textit{Fix: Start the y-axis at 0.}
      \item[Trick 2 --- 3-D or perspective effects]
            Near slices or bars look bigger than far ones even if the data are
            the same.\\
            \textit{Fix: Use a flat 2-D display.}
      \item[Trick 3 --- Inconsistent scale or bar widths]
            Unequal intervals or variable bar widths exaggerate or shrink
            differences.\\
            \textit{Fix: Use uniform intervals and equal bar widths.}
    \end{description}

    \vspace{0.4cm}
    {\color{slate}\footnotesize
      AP prompt: ``Explain one way the graph is misleading and describe how to
      correct it.'' Always name the specific feature and its effect on the reader.}
  \end{minipage}
\end{frame}
}

% ─────────────────────────────────────────────────────────────────────────────
% SLIDE 10 — AP LANGUAGE
% ─────────────────────────────────────────────────────────────────────────────
{
\setbeamercolor{background canvas}{bg=sky}
\begin{frame}[t, plain]
  \navyheader{Interpreting a Bar Chart --- AP Language}
  \vspace{1em}%
  \hspace{0.55cm}%
  \begin{minipage}{0.9\textwidth}
    \small

    \textbf{\color{navy}Every interpretation must include:}
    \begin{itemize}
      \setlength{\itemsep}{0.4em}
      \item The \textbf{total $n$} (``of the 80 adults surveyed\ldots'')
      \item The \textbf{variable} in context (``preferred TV genre'')
      \item The \textbf{most} and \textbf{least} common categories with relative frequencies
      \item A comment on the overall \textbf{shape} of the distribution if relevant
    \end{itemize}

    \vspace{0.5cm}
    {\color{navy}\bfseries\small Model sentence frame:}\\[6pt]
    \small\textit{``Of the \underline{\hspace{1cm}} [group], the most common
    [variable] was \underline{\hspace{1.8cm}} (\underline{\hspace{1cm}}\%).
    The least common was \underline{\hspace{1.8cm}} (\underline{\hspace{1cm}}\%).
    [Additional observation].''}

    \vspace{0.4cm}
    {\color{slate}\footnotesize
      Writing ``40\%'' without context earns \textbf{no credit} on the AP exam.
      Always anchor your numbers to the variable and the group.}
  \end{minipage}
\end{frame}
}

% ─────────────────────────────────────────────────────────────────────────────
% SLIDE 11 — GROUP ACTIVITY
% ─────────────────────────────────────────────────────────────────────────────
{
\setbeamercolor{background canvas}{bg=redbg}
\begin{frame}[t, plain]
  \begin{tikzpicture}[remember picture, overlay]
    \fill[redacc] (current page.north west)
      rectangle ([yshift=-1.35cm]current page.north east);
    \node[anchor=west, text=white, font=\bfseries\large,
          inner xsep=0.55cm, inner ysep=0pt]
      at ([yshift=-0.68cm]current page.north west)
      {Group Activity --- Streaming Service Survey};
  \end{tikzpicture}
  \vspace{1.2cm}\par
  \hspace{0.55cm}%
  \begin{minipage}{0.9\textwidth}
    \small

    \textbf{Data ($n = 50$ students):}
    Netflix~21, Hulu~9, Disney+~10, HBO Max~6, Other~4.

    \vspace{0.35cm}
    \textbf{Part 1 (8 min):}
    \begin{enumerate}
      \setlength{\itemsep}{0.3em}
      \item Complete the relative frequency column.
      \item Construct a correctly labeled relative frequency bar chart.
      \item Write a 2-sentence interpretation.
    \end{enumerate}

    \vspace{0.35cm}
    \textbf{Part 2 (5 min):}
    \begin{enumerate}
      \setlength{\itemsep}{0.3em}
      \item Identify one change to your bar chart that would make it misleading.
            Describe what a reader would incorrectly conclude.
      \item (Optional) Compute sector angles and sketch a pie chart.
            Is the bar chart or pie chart easier to interpret here?
    \end{enumerate}
  \end{minipage}
\end{frame}
}

% ─────────────────────────────────────────────────────────────────────────────
% SLIDE 12 — EXIT TICKET
% ─────────────────────────────────────────────────────────────────────────────
{
\setbeamercolor{background canvas}{bg=sky}
\begin{frame}[t, plain]
  \navyheader{Exit Ticket --- 5 Minutes}
  \vspace{1em}%
  \hspace{0.55cm}%
  \begin{minipage}{0.9\textwidth}
    \small

    A survey of \textbf{80 adults} asks: ``What is your favorite TV genre?''

    \vspace{0.3cm}
    \renewcommand{\arraystretch}{1.5}
    \begin{tabular}{@{} >{\bfseries}p{0.18\textwidth}
                        >{\centering\arraybackslash}p{0.16\textwidth}
                        >{\centering\arraybackslash}p{0.22\textwidth} @{}}
      \rowcolor{navy}
      {\color{white}Genre} & {\color{white}Count} & {\color{white}Rel.\ Freq.} \\
      \midrule
      News    & 32 & ? \\
      Sports  & 20 & ? \\
      Reality & 16 & ? \\
      Other   & 12 & ? \\
    \end{tabular}

    \vspace{0.4cm}
    \begin{enumerate}
      \setlength{\itemsep}{0.35em}
      \item Complete the relative frequency column.
      \item Write one sentence describing the distribution in context.
      \item Describe one way someone could make this bar chart misleading.
    \end{enumerate}
  \end{minipage}
\end{frame}
}

% ─────────────────────────────────────────────────────────────────────────────
% SLIDE 13 — BIG IDEAS / PREVIEW
% ─────────────────────────────────────────────────────────────────────────────
{
\setbeamercolor{background canvas}{bg=navy}
\begin{frame}[t, plain]
  \begin{tikzpicture}[remember picture, overlay]
    \fill[navylight] (current page.north west)
      rectangle ([yshift=-1.35cm]current page.north east);
    \node[anchor=west, text=white, font=\bfseries\large,
          inner xsep=0.55cm, inner ysep=0pt]
      at ([yshift=-0.68cm]current page.north west)
      {Big Ideas \& Preview};
  \end{tikzpicture}
  \vspace{1.2cm}\par
  \hspace{0.55cm}%
  \begin{minipage}{0.9\textwidth}
    \small

    {\color{goldacc}\bfseries Today's core takeaway:}\\[6pt]
    {\color{white}
    Bar charts and pie charts display categorical distributions visually.
    A graph must have a title, labeled axes, and a y-axis starting at 0
    to be truthful.}

    \vspace{0.55cm}
    {\color{skymid}\bfseries Connections:}
    \begin{itemize}
      \setlength{\itemsep}{0.35em}
      \item {\color{white}
            \textbf{Lesson 1.5:} Histograms display \textit{quantitative} data.
            Bars touch because the scale is continuous --- the key difference from today.}
      \item {\color{white}
            \textbf{Unit 2:} Segmented bar charts extend today's displays to
            two categorical variables (two-way tables).}
      \item {\color{white}
            \textbf{AP Exam:} Free-response often asks you to \textit{construct}
            a graph \textit{and} write a complete interpretation in context.}
    \end{itemize}

    \vspace{0.55cm}
    {\color{skymid}\small\itshape
      Complete your homework (Lesson 1.4 HW) before next class.}
  \end{minipage}
\end{frame}
}

\end{document}
